\documentclass[../main.tex]{subfiles}

\begin{document}
	
This chapter evaluates the effectiveness of the proposed topology-aware FlexiCubes framework. The evaluation focuses on determining whether the proposed optimization preserves geometric fidelity while improving the directional consistency of the generated mesh. First, the experimental settings and evaluation metrics are introduced. The experimental methodology is then described, including the datasets, baseline methods, and evaluation procedures. Finally, quantitative and qualitative results are presented and discussed.

\section{Evaluation Parameters}

The proposed method is evaluated using a collection of ShapeNet models belonging to the lamp category. The dataset contains diverse geometric structures with varying levels of complexity, making it suitable for evaluating the robustness of topology-aware mesh extraction.

For each model, an implicit surface is reconstructed using both the original FlexiCubes algorithm and the proposed topology-aware optimization framework. To evaluate the behavior of the generated meshes under different levels of geometric simplification, four target mesh resolutions are considered:
\[
40\,000,\;
20\,000,\;
8\,000,\;
2\,000
\]
triangles.

Mesh simplification is performed using the Quadric Error Metric (QEM) implementation provided by Open3D.

Three quantitative metrics are employed throughout the experiments.

\begin{itemize}
	
	\item Root Mean Square (RMS) surface error, which measures the average geometric deviation between the original mesh and its simplified version.
	
	\item Hausdorff distance, which evaluates the maximum geometric deviation and captures localized reconstruction errors.
	
	\item Normal deviation, computed as the average angular difference between corresponding surface normals, which measures the preservation of local surface orientation.
	
\end{itemize}

Since the objective of the proposed method is to improve topology while maintaining geometry, these complementary metrics provide a comprehensive assessment of reconstruction quality.

%%%%%%%%%%%%%%%%%%%%%%%%%%%%%%%%%%%%%%%%%%%%%%%%%%%%%%%%%%%%%%

\section{Experimental Methodology}

The proposed framework is compared directly with the original FlexiCubes algorithm, which serves as the baseline throughout all experiments. FlexiCubes was selected because it represents the current state-of-the-art differentiable isosurface extraction method and forms the basis of the proposed topology-aware extension.

For each ShapeNet model, the following procedure is performed:

\begin{enumerate}
	
	\item Generate the initial mesh using the original FlexiCubes framework.
	
	\item Generate a second mesh using the proposed topology-aware optimization.
	
	\item Simplify both meshes independently to four predefined triangle counts using quadric error decimation.
	
	\item Sample dense point clouds from both meshes.
	
	\item Compute RMS error, Hausdorff distance, and normal deviation for every level of detail.
	
\end{enumerate}

Ten representative models are evaluated. The reported results correspond to the average performance across all models in order to reduce the influence of individual object characteristics.

All experiments are conducted using identical simplification parameters and evaluation settings for both methods to ensure a fair comparison.

%%%%%%%%%%%%%%%%%%%%%%%%%%%%%%%%%%%%%%%%%%%%%%%%%%%%%%%%%%%%%%
\section{Geometric Accuracy Evaluation}

Tables~\ref{tab:rms_results} and~\ref{tab:hausdorff_results} summarize the geometric reconstruction accuracy of the original FlexiCubes algorithm and the proposed topology-aware optimization under different levels of mesh simplification.

\begin{table}[H]
	\centering
	\caption{Average RMS surface error}
	\label{tab:rms_results}
	\begin{tabular}{lcccc}
		\hline
		& 40k & 20k & 8k & 2k \\
		\hline
		Original Method & 0.001969 & 0.001987 & 0.002050 & 0.002832 \\
		Proposed Method & \textbf{0.001954} & \textbf{0.001961} & \textbf{0.001977} & \textbf{0.002438} \\
		\hline
	\end{tabular}
\end{table}

The proposed method consistently produces lower RMS reconstruction error than the original FlexiCubes algorithm across all levels of detail. The improvement is relatively modest for high-resolution meshes, where both methods closely approximate the original geometry. However, the performance gap becomes more pronounced as the number of triangles decreases. At the most aggressive simplification level of 2,000 triangles, the RMS error decreases from 0.002832 to 0.002438, corresponding to an improvement of approximately 14\%. These results indicate that the proposed topology-aware optimization preserves the overall surface geometry more effectively during mesh simplification.

\begin{table}[H]
	\centering
	\caption{Average Hausdorff distance}
	\label{tab:hausdorff_results}
	\begin{tabular}{lcccc}
		\hline
		& 40k & 20k & 8k & 2k \\
		\hline
		Original Method & 0.029242 & 0.032577 & 0.024871 & 0.032712 \\
		Proposed Method & \textbf{0.006862} & \textbf{0.006805} & \textbf{0.008100} & \textbf{0.027895} \\
		\hline
	\end{tabular}
\end{table}

A more significant improvement is observed in the Hausdorff distance, which measures the maximum geometric deviation between the reconstructed and reference meshes. The proposed method consistently achieves lower Hausdorff distance at every simplification level. The largest improvement occurs for meshes simplified to 20,000 triangles, where the Hausdorff distance decreases from 0.032577 to 0.006805, representing a reduction of approximately 79\%. Even at the lowest resolution of 2,000 triangles, the proposed method maintains a lower maximum geometric deviation than the original algorithm. This demonstrates that the additional topology-aware optimization does not introduce localized geometric artifacts and instead improves geometric robustness throughout the simplification process.

Overall, both the RMS and Hausdorff distance results demonstrate that incorporating topology-aware directional constraints into the FlexiCubes optimization process does not sacrifice reconstruction accuracy. On the contrary, the proposed method consistently produces meshes that are geometrically closer to the original models while simultaneously encouraging more structured mesh topology.

%%%%%%%%%%%%%%%%%%%%%%%%%%%%%%%%%%%%%%%%%%%%%%%%%%%%%%%%%%%%%%
\section{Surface Normal Preservation}

Table~\ref{tab:normal_results} presents the average surface normal deviation measured after mesh simplification.

\begin{table}[H]
	\centering
	\caption{Average normal deviation (degrees)}
	\label{tab:normal_results}
	\begin{tabular}{lcccc}
		\hline
		& 40k & 20k & 8k & 2k \\
		\hline
		Original Method & 24.63 & 24.55 & 24.73 & 26.78 \\
		Proposed Method & 23.97 & 23.95 & 24.49 & 27.44 \\
		\hline
	\end{tabular}
\end{table}

The proposed method achieves slightly lower normal deviation than the original FlexiCubes algorithm for the 40k, 20k, and 8k triangle meshes. The improvements are relatively small, ranging from approximately $0.2^\circ$ to $0.7^\circ$, indicating that both methods preserve surface orientation with comparable accuracy at higher levels of detail. At the lowest resolution of 2,000 triangles, the proposed method exhibits a slightly higher normal deviation than the original method, suggesting that aggressive simplification introduces additional local orientation changes.

Overall, the results demonstrate that incorporating topology-aware optimization has only a minor influence on surface normal preservation. This observation is expected because the proposed method is designed primarily to improve mesh topology and edge-flow organization rather than directly optimize surface normals. More importantly, the additional topology constraints do not substantially degrade the geometric characteristics of the reconstructed meshes.

%\section{Edge Direction Preservation}



\section{Visual Comparison}

\begin{table}[H]
	\centering
	\renewcommand{\arraystretch}{1.2}
	
	\begin{tabular}{c c c}
		\hline
		& \textbf{Original} & \textbf{Proposed Method} \\
		\hline
		
		\includegraphics[width=0.28\textwidth]{Figure/lamp0.jpg} &
		\includegraphics[width=0.28\textwidth]{Figure/lamp0-original-full_cropped.png} &
		\includegraphics[width=0.28\textwidth]{Figure/lamp0-modified-full_cropped.png} \\
		
		\includegraphics[width=0.28\textwidth]{Figure/lamp1.jpg} &
		\includegraphics[width=0.28\textwidth]{Figure/lamp1-original-full_cropped.png} &
		\includegraphics[width=0.28\textwidth]{Figure/lamp1-modified-full_cropped.png} \\
		
		\includegraphics[width=0.28\textwidth]{Figure/lamp3.jpg} &
		\includegraphics[width=0.28\textwidth]{Figure/lamp3-original-full_cropped.png} &
		\includegraphics[width=0.28\textwidth]{Figure/lamp3-modified-full_cropped.png} \\
		
		\hline
	\end{tabular}
	
	\caption{Visual comparison between the original FlexiCubes method and the proposed topology-aware method.}
	\label{tab:visual_compare}
	
\end{table}

Besides quantitative evaluation, visual inspection provides an intuitive assessment of the quality of the generated meshes. Table~\ref{tab:visual_compare} compares the reconstruction results produced by the original FlexiCubes algorithm and the proposed topology-aware method on three representative lamp models. Each row presents the reference image alongside the meshes generated by the two approaches.

For the simple floor lamp (first row), both methods successfully reconstruct the overall geometry. However, the proposed method generates a noticeably smoother and more continuous supporting pole. The surface exhibits fewer irregular triangles, resulting in a cleaner edge flow while preserving the overall shape of the object.

The second example demonstrates a more articulated desk lamp containing multiple joints and thin structural components. The original FlexiCubes output contains numerous irregular triangles around the articulated arm and connection regions, producing a noisy surface appearance. In contrast, the proposed method significantly reduces these artifacts and generates more coherent polygon distributions. The lamp head and supporting arm appear smoother while maintaining the overall geometric structure of the original object.

The third example contains several curved components and thin cylindrical structures. Similar to the previous examples, the proposed method produces a more regular mesh with improved local connectivity, particularly around the curved neck and lamp shade. The supporting cable and base also exhibit fewer disconnected or highly distorted triangles than the original method.

Overall, the visual comparison indicates that the proposed topology-aware optimization improves mesh quality without noticeably altering the reconstructed geometry. Across all three examples, the generated meshes contain more uniform triangle distributions and smoother surface connectivity while preserving the overall object shape. These qualitative observations are consistent with the quantitative evaluation presented in the previous sections, where the proposed method achieved lower reconstruction errors and improved normal consistency. The results suggest that incorporating topology priors into the FlexiCubes extraction process enhances the structural quality of the generated meshes while maintaining high geometric fidelity.




\section{Chapter Conclusion}

Across all evaluated levels of detail, the proposed topology-aware optimization consistently produces meshes with lower RMS error and lower Hausdorff distance than the original FlexiCubes algorithm. The improvements become more pronounced as the target triangle count decreases, suggesting that the proposed directional regularization generates mesh structures that are more robust to simplification.

Although the current normal deviation metric does not reveal a meaningful difference between the two methods, visual inspection of the reconstructed meshes demonstrates that the proposed approach produces more coherent edge directions and smoother edge flow while preserving overall geometric accuracy. These observations support the hypothesis that incorporating learned directional information into the FlexiCubes optimization process can improve mesh quality without sacrificing reconstruction fidelity.

	
\end{document}